\Askhsh{14782}{1}
\begin{ylh}{2}
    \item 3.3. Εξισώσεις 2ου Βαθμού
\end{ylh}
\begin{alist}
\item Να χαρακτηρίσετε τις προτάσεις που ακολουθούν γράφοντας στην κόλλα σας τη λέξη Σωστό ή Λάθος δίπλα στο γράμμα που αντιστοιχεί σε κάθε πρόταση.

\begin{rlist}
\item Τα σημεία $A(x, y)$ και $B(-x, y)$ είναι για κάθε τιμή των $x, y$ συμμετρικά ως προς τον άξονα $y'y$ (όχι $x'x$).
\item Η εξίσωση $x^\nu = a$ με $a < 0$ και $\nu$ περιττό, έχει ακριβώς μία λύση την $-\sqrt[\nu]{|a|}$.
\item Για οποιουσδήποτε πραγματικούς αριθμούς $a, \beta, \gamma, \delta$ ισχύει: Αν $a < \beta$ και $\gamma < \delta$, τότε $a\gamma < \beta\delta$.
\item Για κάθε $\theta \in (0, +\infty)$ ισχύει: $|x| < \theta \Leftrightarrow -\theta < x < \theta$.
\item Αν $x_1, x_2$ οι ρίζες της εξίσωσης $a x^2 + \beta x + \gamma = 0$ με $a \neq 0$ και $\Delta \ge 0$, τότε $S = x_1 + x_2 = -\beta/a$ και $P = x_1 \cdot x_2 = \gamma/a$.
\monades{10}
\end{rlist}
\item Αν $x_1, x_2$ οι ρίζες της εξίσωσης $a x^2 + \beta x + \gamma = 0, a \neq 0$ να δείξετε ότι $P = x_1 \cdot x_2 = \dfrac{\gamma}{a}$. \monades{15}
\end{alist}

